\documentclass[12pt]{article}
\usepackage{amsmath}
\usepackage{amssymb}
\usepackage{mathtools}
\usepackage{listings}
\usepackage{xcolor}

\title{intlang: Formal Semantics and Type System}
\author{}
\date{}

\lstset{
  basicstyle=\ttfamily,
  keywordstyle=\color{blue},
  commentstyle=\color{gray},
  stringstyle=\color{red},
  breaklines=true,
  showstringspaces=false,
}

\newcommand{\judge}[1]{\vdash \, #1}
\newcommand{\env}{\Gamma}
\newcommand{\type}{\tau}

\begin{document}

\maketitle

\section{Lexical Conventions}

\begin{itemize}
  \item \textbf{Identifiers:} An identifier is a letter followed by zero or more letters or digits. Identifiers are case-sensitive.
  \item \textbf{Integer literals:} $ n \in \mathbb{N} \cup \{0\} $. While only non-negative integers can be written down as literals. Negative integers exist in the language and can just be written down as $0 - n$.
  \item \textbf{Whitespace:} Ignored (spaces, tabs, newlines).
\end{itemize}

\section{Concrete Grammar (BNF)}

The concrete grammar is defined in extended Backus-Naur form:

\subsection{Programs and Statements}

\begin{align*}
\text{prog} &::= \text{nlexp}^* \, \text{expr} \\
\text{nlexp} &::= \textbf{let} \, \text{id} \, \textbf{=} \, \text{expr} \, \textbf{;}
\end{align*}

A program consists of zero or more named let-bindings followed by a final expression. Let bindings must appear before any expression.

\subsection{Expressions}

Expressions are defined with explicit precedence levels (from lowest to highest binding strength):

\subsubsection{Lambda}

\begin{align*}
\text{expr} &::= \lambda \, \text{id} \, . \, \text{expr} \\
            &\mid \text{expr}_{\text{cmp}}
\end{align*}

Lambda abstraction binds weakest and is right-associative: $\lambda x. \lambda y. e$ parses as $\lambda x. (\lambda y. e)$.

\subsubsection{Comparison Operators}

\begin{align*}
\text{expr}_{\text{cmp}} &::= \text{expr}_{\text{cmp}} \, \textbf{==} \, \text{expr}_{\text{add}} \\
                         &\mid \text{expr}_{\text{cmp}} \, < \, \text{expr}_{\text{add}} \\
                         &\mid \text{expr}_{\text{add}}
\end{align*}

Comparison operators are left-associative.

\subsubsection{Addition and Subtraction}

\begin{align*}
\text{expr}_{\text{add}} &::= \text{expr}_{\text{add}} \, \textbf{+} \, \text{expr}_{\text{mul}} \\
                         &\mid \text{expr}_{\text{add}} \, \textbf{-} \, \text{expr}_{\text{mul}} \\
                         &\mid \text{expr}_{\text{mul}}
\end{align*}

Addition and subtraction are left-associative.

\subsubsection{Multiplication}

\begin{align*}
\text{expr}_{\text{mul}} &::= \text{expr}_{\text{mul}} \, \textbf{*} \, \text{expr}_{\text{app}} \\
                         &\mid \text{expr}_{\text{app}}
\end{align*}

Multiplication is left-associative.

\subsubsection{Application}

\begin{align*}
\text{expr}_{\text{app}} &::= \text{expr}_{\text{app}} \, \text{expr}_{\text{atom}} \\
                         &\mid \text{expr}_{\text{atom}}
\end{align*}

Application is left-associative and binds strongest: $f \, g \, h$ parses as $(f \, g) \, h$.

\subsubsection{Atomic Expressions}

\begin{align*}
\text{expr}_{\text{atom}} &::= \text{id} \\
                          &\mid n \\
                          &\mid \textbf{(} \, \text{expr} \, \textbf{)}
\end{align*}


\end{document}
